Fix a law \(P\) in either displayed class, an arm \(t\), and a point \(x\in\mathcal B_P\).  Put \(R_x=\lVert X-x\rVert_2\).  Symmetry of the triangular kernel and the signed-distance definition give
\[
 K_\triangle(D_{P,x}/h)=K_\triangle(R_x/h).
\]
Moreover, \(K_\triangle(u)\geq1/2\) for \(0\leq u\leq1/2\).  Therefore, whenever \(0<h\leq2r_0\), actual-arm thickness yields
\[
\begin{aligned}
 \Psi^{(0)}_{t,P,x}(h)
 &=h^{-2}\mathbb E_P\!\left[
 K_\triangle(R_x/h)\mathbf1\{T=t\}\right]\\
 &\geq \frac1{2h^2}P\{T=t,\ R_x\leq h/2\}\\
 &\geq \frac1{2h^2}c_-(h/2)^2
 =\frac{c_-}{8}.
\end{aligned}
\]
The thickness hypothesis is uniform in \(P,t,x\).  For the smooth class, \(\mathcal B_{P,\delta}\subseteq\mathcal B_P\), so the same calculation applies on the trimmed boundary.  Thus one may take \(\kappa_0=c_-/8\), with ``sufficiently small'' meaning \(h\leq2r_0\).